% Owned proposed B.142. All checkpoint29 sources remain frozen.
\subsection*{An actual first-order obstruction to every coarse section}

\statement{B.142}{Proved}{The actual eight-family coarse filling has no section in either category; its one-color first-jet extension class has infinite order and every odd-prime reduction remains nonzero in full completion}
\label{stmt:B142}
Use the actual groups of B.123 and B.127:
\[
 \begin{gathered}
 F=F(U_1,\ldots,U_8)*T,\qquad T=F(f,g,x),\qquad L=\ker(F\to T),\\
 W=F/N,\qquad J=\ker(W\to T),\qquad
 S=F(m,v),\qquad \Gamma_0=W\rtimes S.
 \end{gathered}
Here $N$ is the entire forty-seed normal closure, with both ordinary
node sheets and every signed source and deck translate. The actual
eight-family filling is
\begin{equation}\label{eq:B142-coarse-map}
 \pi:\Gamma_0\longrightarrow Q_0=T\rtimes S.
\end{equation}
Neither $\pi$ nor its map on full profinite completions has a
group-homomorphic right inverse, continuous in the latter category.
No equivariance, prescribed branch or cone lift, common conjugator,
or source normalization is imposed on a proposed section.

More precisely, put $M=J_{\mathrm{ab}}=\mathbb Z[T]^8$ and let
$I\subset\mathbb Z[T]$ be the augmentation ideal. Keeping only the
original $B$ color, whose $F_{11}$ label is $2$, gives the actual
$Q_0$-module quotient
\[
 V=\mathbb Z[T]/I^2\cong\mathbb Z^4.
\]
The resulting actual abelian extension
\begin{equation}\label{eq:B142-integral-extension}
 1\longrightarrow V\longrightarrow E_V
   \longrightarrow Q_0\longrightarrow1
\end{equation}
has an infinite-order class in $H^2(Q_0,V)$. For every odd prime $p$,
its reduction has kernel $V_p=V/pV$ and middle group $E_p$, with
\begin{equation}\label{eq:B142-finite-extensions}
 \begin{gathered}
 1\longrightarrow V_p\longrightarrow E_p
   \longrightarrow Q_0\longrightarrow1,\\
 1\longrightarrow V_p\longrightarrow\widehat E_p
   \longrightarrow\widehat Q_0\longrightarrow1.
 \end{gathered}
\end{equation}
Both corresponding extension classes are nonzero, in ordinary and
continuous cohomology respectively. The second sequence has its
\emph{actual embedded finite kernel}; no completion of $J$ or
identification of the continuous abelianization of $\overline J$
is assumed. These assertions concern the specified coarse filling
and do not decide the remaining fixed-source equality of B.139.

\begin{proof}
\emph{B.142.1. The actual quotient and every possible lift.}
B.127 proves $N\subseteq[L,L]$ for the whole indicated normal closure.
Consequently
\[
 E=\Gamma_0/[J,J]=(M\rtimes T)\rtimes S
\]
is an actual quotient. The source action on the first parenthesis
retains its full corrections; this expression does not split $E$
over $Q_0$. Use the old retained free section
$\sigma:T\to W$, whose generators are
$(f,g,x)=(f,cfc,\tau c)$.

For an arbitrary proposed section $i$ of \eqref{eq:B142-coarse-map},
write
\[
 i(t)=j_t\sigma(t),\qquad i(s)=j_s s,
 \qquad j_t,j_s\in J,quad t\in T,quad s=m,v.
\]
Set $\eta(t)=[j_t]\in M$ and $\lambda_s=[j_s]\in M$.
The homomorphism property on $T$ gives exactly the left derivation rule
\begin{equation}\label{eq:B142-derivation}
 \eta(tu)=\eta(t)+t\eta(u).
\end{equation}
Its three free-generator values are unrestricted. Let $A_s$ be the
actual action on $M$ induced by $\alpha_s=\operatorname{Ad}(s)$,
let $\phi_s$ be the retained action, and put
\[
 d_s(t)=
 [\alpha_s(\sigma(t))\sigma(\phi_s(t))^{-1}]\in M.
\]
The coarse relation $sts^{-1}=\phi_s(t)$ forces
\begin{equation}\label{eq:B142-section-equation}
 A_s\eta(t)+d_s(t)-(\phi_s(t)-1)\lambda_s
       =\eta(\phi_s(t)).
\end{equation}
Indeed the source correction contributes
$\lambda_s-\phi_s(t)\lambda_s$, with this order and sign.
These equations are necessary for every section, without a
selected-involution or paired-color condition.

\emph{B.142.2. The one-color module and its exact first-order action.}
The other seven colors form an invariant submodule: each signed source
preserves its positive peripheral core. On the $B$ color its action is
\[
 \lambda\longmapsto\phi_s(\lambda)h_s.
\]
Every $\phi_s$ preserves $I^2$, and multiplication by the group-ring
unit $h_s$ preserves this two-sided ideal. Retained left multiplication
does so as well. Thus deleting the other colors and $I^2$ is an actual
$Q_0$-module quotient, invariant under all signed source actions.
No selected-involution action is claimed to survive the deletion of
the opposite color.

For $\lambda\in\mathbb Z[T]$, define
\[
 \epsilon(\lambda)=\sum_t\lambda_t,
 \qquad D_r(\lambda)=\sum_t\lambda_t\chi_r(t)
 \quad(r=f,g,x),
\]
where $\chi_r$ is the corresponding exponent character. These give
the coordinates of $V$: the images of
$1,f-1,g-1,x-1$ are an integral basis, using
$uv-1=(u-1)+(v-1)+(u-1)(v-1)$ and the free basis of $T_{\rm ab}$.
For a first-order vector $z=(z_0,z_1,z_2,z_3)$, retained translation is
\begin{equation}\label{eq:B142-retained-jet}
 t z=(z_0,z_1+\chi_f(t)z_0,
              z_2+\chi_g(t)z_0,z_3+\chi_x(t)z_0).
\end{equation}
Both positive source exponent matrices send
$f\mapsto g$, $g\mapsto f$, and $x\mapsto x+f-g$.
Their exact $B$ heads are $h_m=f^{-1}g$ and
$h_v=xg^{-1}x^{-1}$. Hence
\begin{equation}\label{eq:B142-source-jets}
 \begin{aligned}
 A_m z&=(z_0,z_2+z_3-z_0,z_1-z_3+z_0,z_3),\\
 A_v z&=(z_0,z_2+z_3,z_1-z_3-z_0,z_3),\\
 A_m^{-1}&=A_m,\\
 A_v^{-1}z&=(z_0,z_2+z_3+z_0,z_1-z_3,z_3).
 \end{aligned}
\end{equation}
The inverse-source heads are $gf^{-1}$ and $f$, respectively.
The exact retained words needed for the obstruction are
\begin{equation}\label{eq:B142-retained-words}
 \begin{aligned}
 \phi_m(f)&=f^{-1}gf,& \phi_m(g)&=f,& \phi_m(x)&=xg^{-1}f,\\
 \phi_v(f)&=xgx^{-1},&&& \phi_v(x)&=xg^{-1}f.
 \end{aligned}
\end{equation}
The entire relevant components of the actual correction table are
\begin{equation}\label{eq:B142-defects}
 \begin{gathered}
 (d_m(f))_B=(d_m(g))_B=(d_m(x))_B=0,\\
 (d_v(f))_B=1-xg^{-1}x^{-1},\qquad
 (d_v(x))_B=xg^{-1}-xg^{-1}f.
 \end{gathered}
\end{equation}
In particular their required nonzero moments are $1$ and $-1$.
These are complete color identities from the actual marked input,
not selected terms from an unchecked expansion.

For completeness, all signed corrections in the quotient $V$, in
the column order $f,g,x$, are
\[
 \begin{array}{c|ccc}
 s&d_s(f)&d_s(g)&d_s(x)\\ \hline
 m,m^{-1}&(0,0,0,0)&(0,0,0,0)&(0,0,0,0)\\
 v,v^{-1}&(0,0,1,0)&(0,-1,0,0)&(0,-1,0,0).
 \end{array}
\]
Equal first-order vectors do not identify the full signed retained words.
All these actions and corrections descend from the actual group.

\emph{B.142.3. Fourteen unrestricted coordinates and eight equations.}
Write $F_0,\ldots,F_3$, $G_0,\ldots,G_3$, and $X_0,\ldots,X_3$
for the coordinates of $\eta(f)_B,\eta(g)_B,\eta(x)_B$, and put
$M_0=\epsilon((\lambda_m)_B)$,
$V_0=\epsilon((\lambda_v)_B)$.
Applying \eqref{eq:B142-derivation}--\eqref{eq:B142-defects} to
\eqref{eq:B142-section-equation} gives the following necessary equations.
The labels record source, retained generator, and coordinate:
\begin{equation}\label{eq:B142-eight-rows}
 \begin{aligned}
 E_1\ (m,f,0):\quad F_0-G_0&=0,\\
 E_2\ (m,f,1):\quad F_2+F_3-F_0-G_1+G_0&=0,\\
 E_3\ (m,f,3):\quad F_3-G_3&=0,\\
 E_4\ (m,g,1):\quad G_2+G_3-G_0-F_1-M_0&=0,\\
 E_5\ (m,x,1):\quad X_2+X_3-X_0-X_1+G_1-F_1-M_0&=0,\\
 E_6\ (m,x,2):\quad X_1-X_3+X_0-X_2+G_2-F_2+F_0-G_0+M_0&=0,\\
 E_7\ (v,f,2):\quad F_1-F_3-F_0-G_2+X_0-V_0&=-1,\\
 E_8\ (v,x,1):\quad X_2+X_3-X_1+G_1-F_1-V_0&=1.
 \end{aligned}
\end{equation}
For example $\eta(f^{-1}gf)$ has moments
$(G_1-G_0,G_2+F_0,G_3)$, whereas $A_m\eta(f)$ has moments
$(F_2+F_3-F_0,F_1-F_3+F_0,F_3)$. Also
$\eta(xgx^{-1})$ has moments $(G_1,G_2-X_0,G_3+G_0)$.
These calculations fix the derivative convention and the placement
of the peripheral head before any simplification.

The exact combination
\begin{equation}\label{eq:B142-zero-two}
 -E_1-2E_2+E_3+E_4-3E_5-2E_6-E_7+E_8
 \qquad\text{gives}\qquad 0=2.
\end{equation}
Every coefficient of all fourteen unrestricted variables cancels.
One can also eliminate by hand: $E_1$--$E_3$ give
$G_0=F_0$, $G_3=F_3$, and $G_1=F_2+F_3$; adding $E_5,E_6$
gives $F_1=G_2+F_3$, and $E_4$ gives $M_0=-F_0$.
Then $E_5,E_8$ force $V_0=X_0-F_0-1$, while $E_7$ forces
$V_0=X_0-F_0+1$.
Thus no section exists even at the module level. The proof uses only
positive $m,v$ equations and no selected-involution, paired-color,
inverse-source, or bounded-support restriction.

\emph{B.142.4. The actual extension classes.}
The invariant quotient in part~2 gives
$P_V=V\rtimes T$ and $E_V=P_V\rtimes S$, with its actual corrected
source action. Its kernel over $Q_0$ is exactly $V$, proving
\eqref{eq:B142-integral-extension}. In kernel--retained--source
coordinates its product is
\[
 (\mu,t,s)(\nu,u,r)
   =(\mu+tA_s\nu+t d_s(u),\ t\phi_s(u),\ sr).
\]
The canonical set section therefore has extension cocycle
\[
 c((t,s),(u,r))=t d_s(u),
\]
with coefficient action $(t,s)\cdot\mu=tA_s\mu$.
Associativity of the actual group gives the cocycle identity.
Every change of free lifts over $T$ and $S$ is precisely the
derivation/source-correction change used above.

The $n$-fold Baer sum has the same coefficient action and its correction
table multiplied by $n$, by adding the kernel coordinates of the
canonical free lifts in each factor. A split $n$-fold sum would therefore
satisfy \eqref{eq:B142-eight-rows} with both nonzero right sides
multiplied by $n$. Equation \eqref{eq:B142-zero-two} would give $0=2n$.
This is impossible for every nonzero integer $n$, so the class of
\eqref{eq:B142-integral-extension} has infinite order.

For any odd prime $p$, quotienting by $pV$ gives the first sequence
in \eqref{eq:B142-finite-extensions}. The same eight equations over
$\mathbb F_p$ give $0=2$, so its ordinary extension class is nonzero.
This is one uniform certificate for every odd prime, not a collection
of finite-prime tests.

\emph{B.142.5. The finite kernel in the actual full completion.}
Put $P_p=V_p\rtimes T$. The retained action
\eqref{eq:B142-retained-jet} factors through $T_{\rm ab}/pT_{\rm ab}$,
so there is an explicit finite quotient
\[
 P_p\longrightarrow V_p\rtimes(T_{\rm ab}/pT_{\rm ab})
\]
which is injective on $V_p$. Since $P_p$ is finitely generated, its
kernel contains a characteristic finite-index subgroup $U$.
The actual source action preserves $U$ and has finite image on $P_p/U$.
Keeping that action in a finite quotient $\overline S$ gives
\[
 E_p\longrightarrow(P_p/U)\rtimes\overline S
\]
still injective on $V_p$. The initial finite head need not itself be
source-invariant. This proves that the actual finite kernel embeds
in $\widehat E_p$. Quotient exactness and $E_p/V_p=Q_0$ now give
the second sequence in \eqref{eq:B142-finite-extensions}, since the
closure of an embedded finite subgroup is itself.

For clarity, the same characteristic-refinement argument gives full
induced topology of $P_p$ in $E_p$, and of $T$ in $Q_0$. Thus
\[
 \widehat E_p=(V_p\rtimes\widehat T)\rtimes\widehat S,
 \qquad \widehat Q_0=\widehat T\rtimes\widehat S,
\]
with the actual corrected source action on the first parenthesis.
In particular there is an actual continuous quotient
$\widehat\Gamma_0\to\widehat E_p$, whose projection is
$\widehat\pi$. None of these identifications completes the infinite
family kernel separately.

The coefficient action on $V_p$ is continuous because its matrix image
is finite. The compact extension admits a continuous set section:
choose a finite quotient injective on $V_p$, choose representatives
of its quotient by the image of $V_p$, and take the union of their
clopen preimages. Its projection onto $\widehat Q_0$ is a continuous
compact bijection and therefore a homeomorphism. This is only a
set-theoretic section; it defines the continuous extension class.

If this compact extension had a continuous group section, its images
of the three retained generators and two sources would give fourteen
arbitrary $\mathbb F_p$ coordinates satisfying
\eqref{eq:B142-eight-rows}. The same contradiction proves that its
continuous class is nonzero. Finally a section of $\widehat\pi$
would compose with $\widehat\Gamma_0\to\widehat E_p$ to give
such a section. This proves the completed nonsplitting claim, without
residual-finiteness assumptions on $\Gamma_0$ or a continuous
point-coefficient functional on a completed group ring.
\end{proof}

The complete actual coefficient input, the all-signed first-order discovery
system, the separate eight-equation reconstruction without a Smith solver,
and the independent integer certificate are retained at the following paths.
The proof uses the displayed unrestricted certificate, not a search for lifts.
\begin{quote}\footnotesize\raggedright
\path{records/root_coarse_corrections29.json},
\path{work_CE/next_symmetric_nodes/family_kernel_schreier25.json},
\path{work_CE/code/check_arbitrary_equivariant_section_jet29.py},
\path{work_CE/code/check_unrestricted_section_eight_equations29.py},
\path{code/check_root_arbitrary_section_certificate29.py},
\path{work_CE/arbitrary_coarse_section_obstruction29.txt},
\path{work_CE/actual_first_jet_extension_class29.txt}.
\end{quote}
